% =============================================================================
% 04_THEMATIC_APPLICATION — Deep Analysis of Four Cornerstone Papers
% =============================================================================
% Page target: 3–4 pages
% Structure: Introduction → Pham → Tian → Hussein → Alsaigh → Cross-cutting
% =============================================================================

\section{Thematic Application: Four Cornerstone Papers}
\label{sec:thematic}
\addcontentsline{toc}{section}{Thematic Application}

The following four papers form the analytical foundation of this review.
Each represents a distinct but complementary contribution to the intersection
of health AI evaluation and governance, and together they constitute the
primary evidence base from which the review's overarching arguments are
constructed. Individually, they address evaluation methodologies, governance
capacity, organizational maturity, and operational audit frameworks.
Collectively, they reveal a coherent structural logic: from diagnosis of the
problem, to diagnostic tools, to operational instruments.

\subsection{Pham (2025): Ethical and Legal Framework for Health AI}
\label{sec:pham}

\citet{pham2025ethical} provide the broadest framing contribution among the
four cornerstone papers. Published in the Royal Society Open Science, their
review synthesizes ethical and legal considerations in healthcare AI across
four dimensions: innovation, safety, fairness, and regulatory governance.
Unlike papers that address a single facet of the problem, Pham's work maps
the entire terrain — connecting the technical properties of AI systems to
their societal implications and the regulatory architectures that govern them.

Pham's central argument is that the deployment of AI in clinical settings
raises a set of \textit{distinctive} ethical challenges that existing
regulatory frameworks, designed for pharmaceuticals or conventional medical
devices, are ill-equipped to address. These challenges include: the dynamic
nature of continuously learning AI systems, which complicates the assumption
of a fixed risk profile; the opacity of deep learning models, which defeats
traditional disclosure-based regulatory strategies; the scalability of AI
deployment, which can propagate biases at a population scale before
individual clinicians are aware of the problem; and the distribution of
accountability across developers, clinicians, health systems, and regulators,
which no existing legal framework clearly resolves.

The paper's primary contribution to this review is its integrative scope.
Pham demonstrates that evaluation, governance, and implementation are not
separate policy domains but interdependent dimensions of a single challenge.
A validation methodology that ignores the conditions of real-world
deployment will generate misleading performance estimates; a governance
framework that ignores the organizational dynamics of implementation will
fail even when its technical standards are met; and an implementation
strategy that ignores the ethical stakes will erode the trust necessary
for sustainable adoption. This integrative logic is the intellectual
backbone of the present review.

\subsection{Tian et al. (2026): The Pilot Trap and Institutional Capacity}
\label{sec:tian}

If Pham sets the broad framing, \citet{tian2026pilottrap} provide the
empirical diagnosis. Published in the Journal of Medical Internet Research,
their study is among the most methodologically rigorous accounts of health
AI implementation available: an 18-month longitudinal qualitative case
study of a provincial clinical AI platform in China, drawing on observation,
interviews, and operational data from deployment start through routine use.

The paper's conceptual centrepiece is the \textbf{pilot trap}: the observation
that AI systems frequently demonstrate clinical value at the pilot stage but
systematically fail to transition to routine care. The conventional explanation
for this pattern — that the AI is insufficiently accurate or that clinicians
resist technology — is rejected by Tian and colleagues on the basis of their
fieldwork. In their study, the pilot phase was not the problem; the AI
performed adequately within the controlled conditions of the pilot. The
failure occurred precisely at the point of transition, driven by governance
gaps: unclear accountability boundaries between the clinical team and the
AI system; absent coordination mechanisms between departments; inadequate
infrastructure for model updates; and organizational structures that had
no role designated for managing AI lifecycle governance.

Tian and colleagues' most important theoretical contribution is the
\textbf{institutional capacity argument}: governance capacity is not a
precondition that health systems possess or lack, but a capability that
is \textit{built through the act of implementation itself}. This is a
direct empirical challenge to the assumption shared by the FDA and the
EU AI Act — that governance can be specified in advance and verified
prior to deployment. The six-module governance framework they propose
— institutional carrier formation, infrastructure governance, regulatory
and ethical governance, interdisciplinary coordination, translational
scaling, and lifecycle evaluation — was derived empirically from
operational tensions observed during deployment. The modules are not
prescriptive boxes to check before going live; they are analytical
descriptions of what effective governance \textit{looks like when it
emerges}.

For this review, Tian's paper provides the evidence base for the claim
that the governance deficit in health AI is not primarily a funding or
technology problem, but a process problem: governance must be designed
to evolve alongside the AI system, not specified once at the outset.

\subsection{Hussein et al. (2026): A Governance Maturity Model}
\label{sec:hussein}

While Tian's contribution is diagnostic — identifying what goes wrong
and why — \citet{hussein2026governance} offer a prescriptive response.
Published in \textit{NPJ Digital Medicine}, their paper presents a
\textbf{comprehensive maturity model for AI governance in healthcare},
derived from a systematic review of existing frameworks and validated
through expert consultation.

The maturity model proposes five levels of organizational AI governance
maturity, from ad hoc and reactive governance (Level 1: Initial/Ad Hoc)
through to institutionalized, continuously improving governance ecosystems
(Level 5: Leading).
Each level is characterized by specific organizational capabilities:
data governance infrastructure; algorithmic audit and monitoring systems;
clinical governance integration; regulatory compliance mechanisms; and
strategic AI leadership. The model serves two functions: as a diagnostic
tool (enabling organizations to locate themselves on the maturity
spectrum) and as a roadmap (identifying the capabilities that must be
built to advance to the next stage).

The paper's primary relevance to this review lies in its addressing
of the participatory governance theme — the recognition that AI governance
cannot be confined to technical and regulatory experts, and must incorporate
clinical staff, patients, and communities. The maturity model includes
criteria for stakeholder engagement, transparency mechanisms, and
accountability structures, connecting the governance capacity conversation
directly to the participatory governance literature reviewed in
Section~\ref{sec:related_work}.

Methodologically, the maturity model is notable for bridging the gap
between conceptual framework and practical organizational assessment.
Unlike governance taxonomies that describe what \textit{should} exist,
Hussein's model provides measurable criteria for evaluating where
a given organization stands — an important step toward operationalizing
governance evaluation in health systems.

\subsection{Alsaigh et al. (2026): Evidence-Driven Governance via PEARL-PATHWAY}
\label{sec:alsaigh}

The fourth cornerstone paper brings the operational dimension into sharpest
focus. \citet{alsaigh2026pearl} address a specific and consequential
question: how can empirical healthcare evidence be systematically linked to
governance frameworks in a specific implementation context? Their answer —
an evidence-driven approach integrating the PEARL framework with the PATHWAY
analysis method, applied to the healthcare system of Al Madinah, Saudi Arabia
— represents the most contextually grounded governance analysis among the
four papers.

The PEARL framework structures the analysis across four ethical pillars:
\textbf{Proportionality} (governance burdens scaled to risk), \textbf{Equity}
(fair distribution of AI benefits and harms across populations),
\textbf{Accountability} (clear assignment of responsibility across actors),
and \textbf{Legitimacy} (stakeholder endorsement and procedural fairness).
The PATHWAY method then maps these pillars onto the operational realities
of a specific healthcare system: extracting structured representations from
a corpus of 4,277 healthcare publications refined to 243 articles meeting
inclusion criteria, and evaluating the alignment between identified healthcare
priorities and existing governance structures.

The Al Madinah context is methodologically significant. The city's
demographic diversity — including millions of Hajj and Umrah visitors annually
— creates healthcare governance pressures that reveal the limitations of
one-size-fits-all governance frameworks. The paper demonstrates that
adaptive governance must be \textit{context-sensitive}: the same AI system
deployed in different healthcare contexts requires different governance
calibrations depending on patient population characteristics, resource
availability, and institutional capacity. This directly extends Tian's
institutional capacity argument and Hussein's maturity model into operationalized
governance guidance.

The framework's most important innovation is the \textbf{evidence-to-governance
bridge}: rather than proposing governance principles from first principles
or surveying existing frameworks abstractly, Alsaigh and colleagues extract
governance requirements from empirical evidence of actual healthcare needs.
This reverses the usual top-down logic of governance design — instead of
asking ``what should governance achieve?'' they ask ``what does this context
actually require?'' — and in doing so provides a replicable methodology for
context-sensitive governance calibration that can be applied to other
healthcare systems beyond Al Madinah.

\subsection{Cross-Cutting Synthesis}
\label{sec:cross_cutting}

Viewed together, the four cornerstone papers form a coherent intellectual
architecture for understanding and addressing the health AI evaluation
and governance challenge.

Pham provides the broadest frame: situating health AI governance within
the larger challenge of ensuring that technological innovation serves
clinical and societal benefit. Tian provides the empirical diagnosis:
identifying the pilot trap as a governance failure, not a technical
failure, and establishing that governance capacity is built through
implementation. Hussein operationalizes this insight into a diagnostic
tool: a maturity model that allows organizations to locate themselves
on the governance capacity spectrum and to identify the specific gaps
they must close. Alsaigh operationalizes it further into an operational
instrument: a context-sensitive evidence-to-governance bridge with a
replicable methodology for calibrating governance to local healthcare needs.

The logical sequence from framing to diagnosis to diagnostic tool to
operational instrument mirrors the three structural findings of this
review: the ecological validity gap (evaluation failure), the governance
capacity deficit (governance failure), and the pilot trap (implementation
failure). Each paper illuminates a different node of this causal chain,
and the failure to connect them — as separate streams of literature —
is precisely the gap this review addresses.

Furthermore, all four papers converge on the same underlying premise:
that governance in health AI is not a static condition to be achieved
but an \textit{ongoing organizational and institutional process}. This
premise has direct implications for how health AI should be evaluated,
regulated, and implemented — and it is the analytical thread that
connects the remaining sections of this review.